\section{Separate critical and short-residence residuals}
\label{sec:v98-critical}
Keep the original residual $Y=Y_{a,q,r}$ and write $d=2qr$.
We assign the overlap to the short-residence component:
\begin{equation}\label{eq:v98-disjoint-residual}
 S=Y\cap\{\ell\le3d\},\qquad
 C=Y\cap\{\ell>3d\},\qquad
 Y^\varepsilon=S^\varepsilon+C^\varepsilon.
\end{equation}
Here and below a superscript $\varepsilon$ denotes the original
weight $W^\varepsilon$ on the indicated target restriction. The
letters $C,S$ are target sets, not new probability measures.

\begin{lemma}[A full forward interval on the critical component]
\label{lem:v98-critical-interval}
Every target in $C$ satisfies
\begin{equation}\label{eq:v98-critical-support}
 \tau^-\le2d,\qquad 0<z<a,\qquad \tau^+>d.
\end{equation}
For $0\le s\le d$, the trajectory remains in the original exact
event and
\[
 0<z(\Phi_sx)<b,\qquad b=(a^2+2\Lambda_\circ d)^{1/2}.
\]
Moreover some $s\in[0,d]$ has $A(\Phi_sx)<A(x)/2$.
\end{lemma}
\begin{proof}
The complement of the old protected set is
$\{\tau^-\le2d\}\cap(\{z<a\}\cup\{\tau^+\le d\})$.
Since $\ell>3d$, its second alternative is impossible and
$\tau^+>d$. Radial growth proves the upper bound on $z$.
If $A(\Phi_sx)/A(x)\ge1/2$ on this interval, the same-event times
alone contribute $d/(2r)=q$ to the capped clock. This contradicts
$c_r<q$. The direction choice introduces no new comparison-domain
availability assumption.
\end{proof}

\begin{theorem}[Uncut critical loss paid by the first inverse gradient]
\label{thm:v98-critical-certificate}
Assume $b<b_0$. At every finite collision count, for almost every
target roof,
\begin{equation}\label{eq:v98-critical-certificate}
 C^\varepsilon(t)\le C_0\int_t^{t+d}w_b(u)\,du.
\end{equation}
Thus Theorem~\ref{thm:v98-uncut-insertion}, with $L=d$, is an
explicit bound for the complete critical residual, with no positive
lower gradient cutoff. Under \eqref{eq:v98-spectral-capacity},
\begin{equation}\label{eq:v98-critical-height}
 \limsup_m\sup_{R,n,k}m^2\|C^\varepsilon\|_\infty
                              \le C d b^{1/128}.
\end{equation}
Without that additional condition, the same height bound holds for
$C\cap\{z\ge c\}$ at each previously fixed $c>0$. That last
statement does not imply the uncut assertion by exchanging limits.
\end{theorem}
\begin{proof}
For a target selected in Lemma~\ref{lem:v98-critical-interval},
\[
 A(x)/2\le\int_0^d|\partial_uA(\Phi_sx)|\,ds
       \le C\int_0^d A(\Phi_sx)I_m(\Phi_sx)/z(\Phi_sx)\,ds.
\]
Multiply by $W^\varepsilon(x)\le1$. At each fixed $s$ the
same-event transport is injective in the physical transversal; its
image is inside the original exact fiber and has $z<b$. Therefore
summing the right side gives a submeasure of $w_b(t+s)$.
The untruncated measure has already been proved finite at fixed
count. Tonelli and the common kernels prove
\eqref{eq:v98-critical-certificate}. This is a target-height
comparison, not an integral in starting roofs, and no target weight
is replaced by an unweighted terminal flux.

Under \eqref{eq:v98-spectral-capacity}, apply
\eqref{eq:v98-uncut-limit} to the fixed interval of length $d$.
For the assertion with $z(x)\ge c>0$, use radial monotonicity and
integrate \eqref{eq:v94-critical-finite} only above $c$.
The leading layer cake is at most
$C b^{\alpha-1}(d+B_*^{-1})$, and the finite-count error is
$C_{B_*}(1+B_*d)m^2c^{-1/2}\theta_{B_*}^m$.
Fixing $c,d,B_*$ before the collision limit and enlarging $B_*$
afterwards gives the same bound. This proof neither establishes
\eqref{eq:v98-spectral-capacity} nor removes the fixed $c$ inside
that limit.
\end{proof}

For the retained choice $q=1/8$, $a=\sqrt{qr}$, we have $a^2=d/2$.
Consequently \eqref{eq:v98-critical-height}, whenever its stated
hypothesis is verified, becomes
\begin{equation}\label{eq:v98-critical-power}
             \mathcal H_M(C^\varepsilon)\le C(qr)^{257/256}.
\end{equation}
The exponent is $1+(\alpha-1)/2$. It comes from integrating the
first inverse gradient against the critical-volume power, not from
the small mass of $C$ itself. All internal incidences and all initial
positive gradients are present in the finite-count comparison.

\subsection{A finite-count certificate rather than a hidden hypothesis}
The exact cost of the uncut comparison is visible in
\eqref{eq:v98-uncut-insertion}: at each fixed band it contains
\[
 m^2(\mathcal K_m\theta_{B_*}^{2m})^{1/3}.
\]
The theorem estimates the full source with that cost retained. A
source-specific bound on the weighted momentum-fiber multiplicity
would turn it into the vanishing critical-height theorem
\eqref{eq:v98-critical-power}. Finiteness of $N_m$, a pressure bound
for another partition, or a fixed positive-gradient theorem does not
supply that bound. The condition remains an explicit proof obligation.
